% Préambule commun des schémas TikZ — charte Liouaï
% Chaque schéma commence par % Préambule commun des schémas TikZ — charte Liouaï
% Chaque schéma commence par % Préambule commun des schémas TikZ — charte Liouaï
% Chaque schéma commence par \input{_preambule} puis \begin{document}.
\documentclass[tikz,border=6pt]{standalone}
\usepackage[T1]{fontenc}
\usepackage[utf8]{inputenc}
\usepackage[scaled=0.95]{helvet}
\renewcommand{\familydefault}{\sfdefault}
\usepackage{amsmath,amssymb}
\usepackage[european]{circuitikz}
\usetikzlibrary{arrows.meta,positioning,calc,angles,quotes,decorations.pathreplacing,backgrounds}

% Couleurs du logo Liouaï Digital Education
\definecolor{navy}{HTML}{243743}
\definecolor{vert}{HTML}{28B78D}
\definecolor{vertfonce}{HTML}{167A5C}
\definecolor{vertclair}{HTML}{EAF7F2}
\definecolor{gris}{HTML}{A1A9AE}
\definecolor{gristexte}{HTML}{5C6870}
\definecolor{fond}{HTML}{F3F5F6}
\definecolor{jaune}{HTML}{E5C800}
\definecolor{jauneclair}{HTML}{FFFBD6}

\tikzset{
  every picture/.style={line cap=round, line join=round, color=navy, text=navy},
  fleche/.style={-{Stealth[length=2.6mm,width=2mm]}, line width=1.1pt},
  vecteur/.style={fleche, line width=1.4pt},
  titre/.style={font=\bfseries\small, text=navy},
  note/.style={font=\footnotesize, text=gristexte},
  boite/.style={draw=navy, line width=0.9pt, rounded corners=5pt, fill=fond, align=center, inner sep=6pt},
}
 puis \begin{document}.
\documentclass[tikz,border=6pt]{standalone}
\usepackage[T1]{fontenc}
\usepackage[utf8]{inputenc}
\usepackage[scaled=0.95]{helvet}
\renewcommand{\familydefault}{\sfdefault}
\usepackage{amsmath,amssymb}
\usepackage[european]{circuitikz}
\usetikzlibrary{arrows.meta,positioning,calc,angles,quotes,decorations.pathreplacing,backgrounds}

% Couleurs du logo Liouaï Digital Education
\definecolor{navy}{HTML}{243743}
\definecolor{vert}{HTML}{28B78D}
\definecolor{vertfonce}{HTML}{167A5C}
\definecolor{vertclair}{HTML}{EAF7F2}
\definecolor{gris}{HTML}{A1A9AE}
\definecolor{gristexte}{HTML}{5C6870}
\definecolor{fond}{HTML}{F3F5F6}
\definecolor{jaune}{HTML}{E5C800}
\definecolor{jauneclair}{HTML}{FFFBD6}

\tikzset{
  every picture/.style={line cap=round, line join=round, color=navy, text=navy},
  fleche/.style={-{Stealth[length=2.6mm,width=2mm]}, line width=1.1pt},
  vecteur/.style={fleche, line width=1.4pt},
  titre/.style={font=\bfseries\small, text=navy},
  note/.style={font=\footnotesize, text=gristexte},
  boite/.style={draw=navy, line width=0.9pt, rounded corners=5pt, fill=fond, align=center, inner sep=6pt},
}
 puis \begin{document}.
\documentclass[tikz,border=6pt]{standalone}
\usepackage[T1]{fontenc}
\usepackage[utf8]{inputenc}
\usepackage[scaled=0.95]{helvet}
\renewcommand{\familydefault}{\sfdefault}
\usepackage{amsmath,amssymb}
\usepackage[european]{circuitikz}
\usetikzlibrary{arrows.meta,positioning,calc,angles,quotes,decorations.pathreplacing,backgrounds}

% Couleurs du logo Liouaï Digital Education
\definecolor{navy}{HTML}{243743}
\definecolor{vert}{HTML}{28B78D}
\definecolor{vertfonce}{HTML}{167A5C}
\definecolor{vertclair}{HTML}{EAF7F2}
\definecolor{gris}{HTML}{A1A9AE}
\definecolor{gristexte}{HTML}{5C6870}
\definecolor{fond}{HTML}{F3F5F6}
\definecolor{jaune}{HTML}{E5C800}
\definecolor{jauneclair}{HTML}{FFFBD6}

\tikzset{
  every picture/.style={line cap=round, line join=round, color=navy, text=navy},
  fleche/.style={-{Stealth[length=2.6mm,width=2mm]}, line width=1.1pt},
  vecteur/.style={fleche, line width=1.4pt},
  titre/.style={font=\bfseries\small, text=navy},
  note/.style={font=\footnotesize, text=gristexte},
  boite/.style={draw=navy, line width=0.9pt, rounded corners=5pt, fill=fond, align=center, inner sep=6pt},
}

\begin{document}
\begin{tikzpicture}[x=0.38cm, y=0.38cm]
\def\img{{
{0,0,5,13,9,1,0,0},{0,0,13,15,10,15,5,0},{0,3,15,2,0,11,8,0},{0,4,12,0,0,8,8,0},
{0,5,8,0,0,9,8,0},{0,4,11,0,1,12,7,0},{0,2,14,5,10,12,0,0},{0,0,6,13,10,0,0,0}}}
\foreach \i in {0,...,7}{\foreach \j in {0,...,7}{
  \pgfmathsetmacro{\v}{\img[\i][\j]*100/16}
  \fill[navy!\v!white] (\j,-\i) rectangle ++(1,-1);}}
\draw[gris] (0,0) grid[step=1] (8,-8); \draw[navy, line width=0.8pt] (0,0) rectangle (8,-8);
\node[note] at (4,-9) {image 8 $\times$ 8, gris de 0 à 16};
\draw[fleche, navy] (9.2,-4) -- (12.2,-4) node[midway, above, font=\small] {aplatir};
\foreach \k [count=\c from 0] in {0,0,5,13,9,1,0,0,0,0,13,15}{
  \pgfmathsetmacro{\v}{\k*100/16}
  \fill[navy!\v!white] (13+\c,-3.5) rectangle ++(1,-1);}
\draw[gris] (13,-3.5) grid[step=1] (25,-4.5); \node at (26.3,-4) {$\cdots$};
\draw[navy, line width=0.8pt] (13,-3.5) rectangle (25,-4.5);
\node[note, align=center] at (19.5,-6.2) {un vecteur $\mathbf{x}$ de $\mathbb{R}^{64}$ :\\ la ligne 1, puis la ligne 2, \dots};
\node[boite, anchor=west, align=left, font=\small] at (28.5,-4) {deux images se ressemblent\\ si leurs vecteurs sont proches :\\[3pt]
$\lVert \mathbf{x} - \mathbf{y} \rVert$ petite, ou\\ $\cos(\mathbf{x}, \mathbf{y})$ proche de 1};
\end{tikzpicture}
\end{document}
